\section{Discussion}
\label{sec:discussion}

\paragraph{When does memory matter?}
Memory matters when task-relevant state is hidden and must be carried across time to select the right action. In this task family, that state is usually the hidden goal under WeakPO: \task{NoMemory} fails on 7/9 tasks, while either memory-enabled variant recovers. This indicates that the dominant failure is an \emph{information-management} problem before it is a model-capacity problem.

\paragraph{Why does simple persistence suffice here?}
The current tasks primarily require point memory: preserving the hidden goal matters more than how that value is packaged. This is why the paper should be read as a diagnostic result, not as a broad ranking of memory architectures. We expect $\Delta_{\text{struct}}>0$ mainly outside this single-goal regime, for example in tasks involving object identity, relational state, or ordered subgoals.

\paragraph{Failure mode and practical implication.}
WeakPO exposes goal-seeking failure, which is largely resolved by goal memory. StrongPO exposes object-localization failure, where goal-only memory is insufficient and the next intervention should be object-state estimation, object buffers, or belief estimation. Thus, for this benchmark family, the practical order is \emph{cache first, learned recurrent memory second, structured memory third}, with recovery treated separately because retries and rollback consume finite-horizon budget.
